\section{A Sub-Exponential Time Approximate Counter}

In the previous subsections, we had presented various (randomized) polynomial-time algorithms in order to estimate the number of minimal execution graphs of the given concurrent program. Due to the $\sharpP$-hardness result of approximation (Theorem~\ref{thm:approximate-counting-problem-hardness}), all
of these algorithms, in the worst case, are going to have exponentially high variance.
In light of this, we can ask the following question: 
What if we relax the requirement that our estimation algorithm must work in polynomial time? 
%Of course, if we allow the algorithm to work in exponential time, then we can get zero variance, 
%by traversing the entire \trust{} tree completely deterministically. 
%Hence, in order for this question to make sense, we ask whether there is an
%$(r, 1/4)$-approximate counter for constant $r$ that runs in sub-exponential time.
%Surprisingly, we prove that such an estimation algorithm exists. 

More precisely, note that as mentioned in Section~\ref{trust_sec}, the depth and the branching factor of the \trust tree $\mathcal{D}(P)$ is quadratic in $|P|$. Hence, the total size of $\mathcal{D}(P)$ could be $M := O((|P|^2)^{|P|^2+1} - 1)$ in the worst case. In such cases, we would like to know whether we can avoid exhaustive exploration of $\mathcal{D}(P)$ while still getting a good estimate.
Surprisingly, we prove that such an estimation algorithm exists.

%\begin{theorem}
%\label{th:stockmeyer}
 %   For any constants $r$ and $\rho$, there is a randomized $\tilde{O}(\sqrt{C(P)})$-time $(r, \rho)$-approximate counter.
%\end{theorem}

\begin{theorem}
\label{th:stockmeyer}
    For any constants $r$ and $\rho$, there is a randomized $\tilde{O}(\sqrt{M})$-time $(r, \rho)$-approximate counter.
\end{theorem}

The proof can be found in \cref{sec:stockmeyer}.
Here $\tilde{O}$ suppresses terms that are a polynomial in $|P|$. Note that this theorem is quite surprising: 
%First, suppose we consider the case where $C(P)$ is a polynomial in $|P|$. Then the approximate counter promised by this theorem runs in polynomial-time and gives an approximation of $C(P)$ that is at most a constant factor away from the actual value of $C(P)$. This is in stark contrast to the other algorithms that we have seen in this paper, which can all have exponential variance, even when $C(P)$ is a linear function of $|P|$ (as evidenced by the examples \ref{ex:r+r+r}, \ref{ex:r+rr} and \ref{ex:hairbrush}, of which the latter can be easily generalized to produce instances of exponential variance for \textbf{Algorithm S} as well). 
It shows that even when $C(P)$ is quite close to $M$, we can avoid exhaustive exploration and still get a good estimate of $C(P)$. 
%For instance, as we shall see below, \RM{where do we see this?} our algorithm indicates that if $M$ is around 1 million, then with about 1000 samples, we can get a good estimate with high probability!
However, the catch behind this algorithm is that the search has a high polynomial cost, i.e., the polynomial term suppressed by the $\tilde{O}$ notation is roughly $|P|^6$ and so we did not implement it.

We now describe the main ideas behind this algorithm, which go back to Stockmeyer \cite{Stockmeyer85}.
To describe our algorithm, we assume that we are given a concurrent program $P$, 
which defines the tree $\mathcal{D}(P)$. 
As mentioned before, if $d$ is the size of the program $P$, 
this tree has height $d^2$ and the degree of each node is at most $d^2$. 
However, every node need not have degree $d^2$. In order to make this tree uniform, we add dummy nodes to each node of $\mathcal{D}(P)$ so that the degree of each node becomes $d^2$. 
Call the resulting tree $\mathcal{H}(P)$. Note that $\mathcal{H}(P)$ 
is the complete $d^2$-ary tree of height $d^2$. 
(We will never explicitly construct $\mathcal{D}(P)$ nor $\mathcal{H}(P)$ during our algorithm; 
these are just defined implicitly). 

The algorithm distinguishes two situations: either $C(P)$ is ``small'' (less than a parameter $\theta$ that we will select later)
or ``big'' (bigger than the parameter $\theta$).
If $C(P)$ is big, then we independently sample a number of random walks proportional to $M/\theta$ in $\mathcal{H}(P)$.
We then count the number of times we hit a leaf node of $\mathcal{D}(P)$ in these random walks and then return this fraction of samples as our estimate.
Using Chebyshev bounds, we can bound the error in the procedure.
On the other hand, if $C(P)$ is small, we can enumerate all elements by a binary search.
The binary search uses an inorder numbering of all the nodes of the tree.
The numbering is used to define intervals of leaves, and the search checks if there is a leaf of $\mathcal{D}(P)$ in a given interval.
This query becomes the same as checking if a given path in the tree $\mathcal{H}(P)$ can be extended to a maximal execution graph in $\mathcal{D}(P)$, for which we can use \trust.

Finally, we set the parameter $\theta$ to be $\sqrt{M}$ using which we can bound the entire run time by $\tilde{O}(\sqrt{M})$. 
%We provide the details of the construction in \cref{sec:stockmeyer}.

%\Cref{th:stockmeyer} is extremely surprising: we get an estimate for $C(P)$ by making (roughly) $\sqrt{C(P)}$ runs of the model checker!
%Consider a program with about 1M executions; with about 1K runs, we can get a good estimate with high probability!
%The catch is that the search has a high polynomial cost $O(|P|^6)$ and so we did not implement it.
